\section{The arbitrary-defect upper bound}\label{sec:upper}

We now bound the mass of every configuration whose scores are all
nonpositive.  All degrees, paths and distances in this section belong to the
original tree.  In particular, the argument retains the empty branches:
their contributions will be cancelled separately from the contributions of
edges with occupied support on both sides.

Let $T$ be a finite tree with at least two vertices, and let $C$ satisfy
$S_v(C)\leq 0$ for every $v$.  Set
\[
 \delta_v=-S_v(C)\geq 0.
\]
If $C$ is identically zero the desired bound is immediate.  Henceforth choose
an occupied vertex $o$, so that $C(o)>0$.  This vertex is used to assign edge
owners; the estimator root obtained at the end may be different.

\subsection{Retained edges and the local defect equations}

For adjacent vertices $u,v$, write $T_{u\mid v}$ for the component containing
$u$ after deleting $uv$.  Call $uv$ \emph{retained} if both
$T_{u\mid v}$ and $T_{v\mid u}$ contain an occupied vertex.  Its two messages
are then integers.  For any edge, let
\[
 \overline d_{u\to v}=
 \begin{cases}
 d_{u\to v},&T_{u\mid v}\text{ is occupied},\\
 0,&T_{u\mid v}\text{ is empty}.
 \end{cases}
\]
This notation records the numerical contribution of a message in a score;
it does not identify the message $\EMPTY$ with the integer message
zero.  Thus
\begin{equation}\label{eq:upper-score}
 C(v)=-\delta_v-\sum_{u\sim v}\overline d_{u\to v}.
\end{equation}

\begin{lemma}[Support separation]\label{lem:upper-support}
The two sides of an edge partition $V(T)$.  Every nonretained edge has
exactly one empty side.  A retained edge lies on a path between two occupied
vertices, and every edge of a path between two occupied vertices is retained.
More specifically, if $uv$ is retained and $x$ is an occupied vertex in
$T_{u\mid v}$, every edge on the path from $u$ to $x$ is retained.  If
$uv$ is retained and no other retained edge is incident with $u$, then
$C(u)>0$.
Finally, an endpoint on the empty side of an edge is incident with no
retained edge.
\end{lemma}

\begin{proof}
Deleting an edge of a tree produces exactly two components, giving the
partition.  Since $C(o)>0$, at least one component is occupied.  If both
are occupied, choose one occupied vertex in each; their unique path crosses
the deleted edge.  Conversely, deleting any edge on the path between two
occupied vertices separates those vertices, so that edge is retained.

For the more specific assertion, choose an occupied vertex $y$ in
$T_{v\mid u}$.  The path from $x$ to $y$ consists of the path from $x$ to
$u$, the edge $uv$, and the path from $v$ to $y$.  Every edge on its
$x$--$u$ segment is therefore retained.  If $u$ is unoccupied, an occupied
vertex $x$ on its side is distinct from $u$, and the edge of that
segment incident with $u$ is a retained edge other than $uv$.  This proves
the next assertion by contraposition.

For the last assertion, let $T_{v\mid u}$ be empty.  The edge $uv$ itself
is nonretained.  For every other neighbour $w$ of $v$, the component
$T_{w\mid v}$ lies inside $T_{v\mid u}$ and is empty.  Hence $wv$ is
also nonretained.
\end{proof}

For a retained edge $uv$, put $a=d_{u\to v}$ and $b=d_{v\to u}$.
Removing the opposite message from each endpoint score gives the two
effective branch inputs.  The message recursion therefore yields
\begin{equation}\label{eq:upper-edge-equations}
 a=F(-\delta_u-b),\qquad b=F(-\delta_v-a).
\end{equation}
These equations are only asserted for retained edges.

\begin{lemma}[Classification of retained edge states]\label{lem:upper-classification}
The solutions of \eqref{eq:upper-edge-equations} with
$\delta_u,\delta_v\geq0$ have exactly the following forms, up to exchanging
$u$ and $v$:
\begin{enumerate}
 \item An \emph{oriented state}, with
 \begin{equation}\label{eq:upper-oriented}
 a=k\geq0,\qquad b=-2(k+\delta_v)-3,
 \qquad
 \delta_u=2\delta_v\ \text{or}\ \delta_u=2\delta_v+3.
 \end{equation}
 The latter alternative requires $k>0$.  We direct this edge $u\to v$.
 \item A \emph{two-negative state}, with unique $r,q\in\NN$ such that
 \begin{equation}\label{eq:upper-negative}
 a=-(2r+1),\quad b=-(2q+1),\quad
 \delta_u=r+2q,\quad\delta_v=q+2r.
 \end{equation}
\end{enumerate}
Both messages cannot be nonnegative.
\end{lemma}

\begin{proof}
The definition of $F$ gives
\begin{align*}
 F(z)<0&\iff z\leq1,\\
 F(z)=-(2r+1)&\iff z=1-r\qquad(r\geq0),\\
 F(z)=k&\iff
   \bigl(k>0\text{ and }z=2k\bigr)\text{ or }z=2k+3
   \qquad(k\geq0).
\end{align*}
Indeed, the first two statements use $F(z)=2z-3$ for $z\leq1$;
the even positive and odd positive cases of $F$ give the third, including
the unique preimage $3$ of zero.

If $a=k\geq0$, then $-\delta_v-k\leq0$, so the second edge equation
gives $b=-2(k+\delta_v)-3<0$.  Substituting this into the first equation
and using the nonnegative preimages gives either
$-\delta_u+2k+2\delta_v+3=2k+3$, or the same left side equals $2k$
with $k>0$.  These are exactly the two defect alternatives in
\eqref{eq:upper-oriented}.  The case $b\geq0$ is symmetric.

If both messages are negative, the first two identities express them
uniquely as $-(2r+1)$ and $-(2q+1)$, with $r,q\geq0$.
Their unique preimages give
$-\delta_u+2q+1=1-r$ and
$-\delta_v+2r+1=1-q$, proving \eqref{eq:upper-negative}.
Conversely, substitution verifies every displayed state, subject to the
stated restriction when $k=0$.
\end{proof}

Call an oriented state \emph{defective} if its head defect $\delta_v$ is
positive.  Call a two-negative state defective if $r+q>0$.  The remaining
two-negative state is the neutral state $(-1,-1)$, with both endpoint
defects zero.  The defective retained edges form a forest, since they are
a subset of the edges of $T$.

\subsection{Owners, orientations and auxiliary weights}

\begin{lemma}[Injective incident owners]\label{lem:upper-owners}
For each edge $e=uv$, let $\omega(e)$ be the endpoint farther from $o$.
Then $\omega$ is well defined and injective on $E(T)$.  If
$T_{v\mid u}$ is empty, then $\omega(uv)=v$.
\end{lemma}

\begin{proof}
The unique path from $o$ crosses $uv$ precisely when reaching the endpoint
on the side not containing $o$.  Thus the endpoint distances differ by
one.  The owner is the more distant endpoint.  Each vertex other than $o$
has exactly one neighbour preceding it on its unique path from $o$;
consequently it owns exactly that predecessor edge.  In particular two
edges cannot share an owner.  If $T_{v\mid u}$ is empty, the occupied
vertex $o$ lies on the $u$ side, so $v$ is the farther endpoint.
\end{proof}

Define an auxiliary partial orientation of $T$ as follows.  Keep the direction
of every oriented state, including those with zero head defect.  Direct each
defective two-negative edge toward its owner.  Leave neutral edges and
nonretained edges unoriented.  The classification ensures that no edge has
both directions.  Since $T$ has no undirected cycle, this directed graph has
no directed cycle.

Let $h(v)$ be the maximum length of a directed path ending at $v$, allowing
the path of length zero, and set
\begin{equation}\label{eq:upper-weights}
 A_v=2^{h(v)}.
\end{equation}
The maximum exists because a directed path has no repeated vertex.
For an auxiliary arrow $u\to v$, a longest path ending at $u$ can be
extended by this arrow: if it already contained $v$, it would give a
directed cycle.  Hence
\begin{equation}\label{eq:upper-doubling}
 h(v)\geq h(u)+1,\qquad A_v\geq2A_u.
\end{equation}
A vertex has positive height exactly when it has an incoming auxiliary
arrow.  Since every auxiliary arrow is retained, Lemma~\ref{lem:upper-support}
also shows that the endpoint on the empty side of any edge has height zero
and weight one.

\subsection{Local charges and global defect cancellation}

For an undirected edge $e=uv$, define its weighted message contribution by
\[
 E_e=-\bigl(A_v\overline d_{u\to v}
             +A_u\overline d_{v\to u}\bigr).
\]
Define two nonnegative charges, with disjoint conditions of application:
\begin{align*}
 Q_e^{\mathrm{ret}}
  &=\begin{cases}
     A_{\omega(e)}\delta_{\omega(e)},&e\text{ is retained and defective},\\
     0,&\text{otherwise},
    \end{cases}\\
 Q_e^{\mathrm{emp}}
  &=\begin{cases}
     A_{\omega(e)}\delta_{\omega(e)},&e\text{ is nonretained},\\
     0,&\text{otherwise}.
    \end{cases}
\end{align*}

\begin{lemma}[Retained-edge charge]\label{lem:upper-retained-charge}
For every retained edge $uv$,
\[
 E_{uv}\leq A_u+A_v+Q_{uv}^{\mathrm{ret}}.
\]
\end{lemma}

\begin{proof}
First consider an oriented state $u\to v$ with parameter $k\geq0$.
By \eqref{eq:upper-oriented} and \eqref{eq:upper-doubling},
\begin{align*}
 E_{uv}
 &=A_u\bigl(2(k+\delta_v)+3\bigr)-A_vk\\
 &=3A_u+k(2A_u-A_v)+2A_u\delta_v\\
 &\leq A_u+A_v+2A_u\delta_v.
\end{align*}
The excess $2A_u\delta_v$ is zero when the edge is nondefective.
Otherwise either endpoint can pay it: the local defect relation gives
$A_u\delta_u\geq2A_u\delta_v$, and doubling gives
$A_v\delta_v\geq2A_u\delta_v$.  Thus it is bounded by the charge at
the prescribed owner, regardless of whether that owner is the tail or head.

For a defective two-negative state, name its endpoints so that the auxiliary
direction is $u\to v$, with $v=\omega(uv)$, and use the parameters
in \eqref{eq:upper-negative}.  Then
\begin{align*}
 E_{uv}
 &=A_u+A_v+2A_vr+2A_uq\\
 &\leq A_u+A_v+A_v(2r+q)\\
 &=A_u+A_v+A_v\delta_v.
\end{align*}
Here $q\geq0$ and $2A_u\leq A_v$ justify the inequality.  A neutral
edge has $r=q=0$ and contributes exactly $A_u+A_v$, completing the proof.
\end{proof}

\begin{lemma}[Empty-side cancellation]\label{lem:upper-empty-charge}
If $T_{v\mid u}$ is empty, then
\[
 E_{uv}=A_v\delta_v=Q_{uv}^{\mathrm{emp}}.
\]
\end{lemma}

\begin{proof}
The message $m_{v\to u}$ is literally $\EMPTY$.  At $v$,
the initial pile is zero, and every incident branch other than the $u$ side
is empty.  Consequently
$S_v(C)=d_{u\to v}$, where this last message is an integer because the
$u$ side contains $o$.  Thus
$\overline d_{v\to u}=0$ and $d_{u\to v}=-\delta_v$, giving
$E_{uv}=A_v\delta_v$.  By Lemma~\ref{lem:upper-owners}, $v$ is the
owner.  In addition $A_v=1$, as observed above.  No integer defect-state
equations were applied to the empty message.
\end{proof}

\begin{proposition}[Weighted cancellation]\label{prop:upper-weighted}
The auxiliary weighted mass satisfies
\begin{equation}\label{eq:upper-weighted-bound}
 \sum_v A_vC(v)\leq\sum_v\deg_T(v)A_v.
\end{equation}
\end{proposition}

\begin{proof}
Grouping \eqref{eq:upper-score} by undirected edges gives the exact identity
\begin{equation}\label{eq:upper-weighted-identity}
 \sum_v A_vC(v)=\sum_{e\in E(T)}E_e-\sum_v A_v\delta_v.
\end{equation}
Lemmas~\ref{lem:upper-retained-charge} and \ref{lem:upper-empty-charge}
give, uniformly over all ambient edges,
\[
 E_{uv}\leq A_u+A_v+Q_{uv}^{\mathrm{ret}}+Q_{uv}^{\mathrm{emp}}.
\]
Each sum of the two charges is at most its one owner budget
$A_{\omega(e)}\delta_{\omega(e)}$.  Owners are injective on all edges,
and every vertex budget is nonnegative, so
\[
 \sum_e\bigl(Q_e^{\mathrm{ret}}+Q_e^{\mathrm{emp}}\bigr)
 \leq\sum_e A_{\omega(e)}\delta_{\omega(e)}
 \leq\sum_v A_v\delta_v.
\]
Substitution in \eqref{eq:upper-weighted-identity} cancels the budgets.
Finally, each vertex weight occurs once for each incident edge, so
$\sum_{uv\in E(T)}(A_u+A_v)=\sum_v\deg_T(v)A_v$.
\end{proof}

\subsection{Removing leaf weights}

Let $L$ be the number of ambient vertices of degree one.

\begin{lemma}[Ambient leaf slack]\label{lem:upper-leaf-slack}
Every degree-one vertex $v$ with $A_v>1$ is occupied.  Consequently
\begin{equation}\label{eq:upper-compact}
 \mass{C}\leq L+\sum_{\deg_T(v)>1}\deg_T(v)2^{h(v)}.
\end{equation}
\end{lemma}

\begin{proof}
If $A_v>1$, then $h(v)>0$, so there is an incoming auxiliary arrow at
$v$.  Its underlying edge is retained.  As $v$ has degree one, its side
of that edge is the singleton $\{v\}$; retention implies $C(v)\geq1$.
For any degree-one vertex, therefore,
\[
 A_v-1\leq (A_v-1)C(v):
\]
if $A_v=1$ both sides are zero, and otherwise $C(v)\geq1$ proves it.
Since $A_v\geq1$ at every vertex, all configuration slack terms are
nonnegative.  Using Proposition~\ref{prop:upper-weighted},
\begin{align*}
 \mass{C}
 &=\sum_v A_vC(v)-\sum_v(A_v-1)C(v)\\
 &\leq\sum_v\deg_T(v)A_v
       -\sum_{\deg_T(v)=1}(A_v-1)\\
 &=L+\sum_{\deg_T(v)>1}\deg_T(v)A_v.
\end{align*}
The last equality uses that every vertex of a nontrivial tree has positive
degree.  Replacing $A_v$ by $2^{h(v)}$ proves the claim.  Thus occupied
leaves pay the required slack, while unoccupied ambient leaves have unit
weight and require no payment.
\end{proof}

\subsection{Consolidation to one ambient root}

We prove a pointwise statement about heights before applying the estimator.
A set of vertices is called connected here if the unique tree path between
any two of its vertices stays in the set.  This agrees with connectedness
of the induced subgraph, by uniqueness of paths in a tree.

\begin{lemma}[Source partition]\label{lem:upper-source-partition}
The vertices of $T$ can be partitioned into connected sets $P_i$ with
vertices $s_i\in P_i$ such that
\[
 h(v)=\dist_T(s_i,v)\qquad(v\in P_i).
\]
\end{lemma}

\begin{proof}
For every $v$ of positive height, choose once and for all an incoming
auxiliary arrow $p(v)\to v$ with $h(p(v))=h(v)-1$.  Such an arrow
exists by taking the last edge of a longest directed path ending at $v$:
its prefix has length $h(v)-1$, and \eqref{eq:upper-doubling} at the level
of heights prevents the height of its penultimate vertex from being larger.

Iterating $p$ decreases height by exactly one each time and ends at a
vertex $s(v)$ of height zero after $h(v)$ steps.  No vertex repeats, so
this predecessor chain is the unique tree path from $v$ to $s(v)$ and
$h(v)=\dist_T(s(v),v)$.  Every vertex on the chain has the same terminal
source, since its subsequent predecessor choices are fixed.  Each nonempty
fibre $P_s=\{v:s(v)=s\}$ therefore contains $s$ and the full path from
each of its vertices to $s$.  The path between two vertices of $P_s$ is
contained in the union of their paths to $s$, hence stays in $P_s$.
The nonempty fibres give the required connected partition.
\end{proof}

\begin{lemma}[Merging height-dominated parts]\label{lem:upper-merge}
Let $A,B$ be disjoint nonempty connected vertex sets joined by an edge
$xy$, with $x\in A$ and $y\in B$.  Suppose $a\in A$, $b\in B$ and
$g:V(T)\to\NN$ satisfy
\[
 g(v)\leq\dist_T(a,v)\quad(v\in A),\qquad
 g(v)\leq\dist_T(b,v)\quad(v\in B).
\]
Then $A\cup B$ is connected, and one of $a,b$ satisfies the same domination
on all of $A\cup B$.
\end{lemma}

\begin{proof}
Paths within the two sets, together with $xy$, connect their union.
If $z\in A$ and $w\in B$, concatenate the path from $z$ to $x$ in
$A$, the edge $xy$, and the path from $y$ to $w$ in $B$.  Disjointness
makes this a simple path, so it is the unique tree path and
\begin{equation}\label{eq:upper-boundary-distance}
 \dist_T(z,w)=\dist_T(z,x)+1+\dist_T(y,w).
\end{equation}
Put $\alpha=\dist_T(a,x)$ and $\beta=\dist_T(b,y)$.
If $\beta\leq\alpha+1$, then for every $v\in B$ the triangle
inequality and \eqref{eq:upper-boundary-distance} give
\[
 g(v)\leq\dist_T(b,v)
 \leq\beta+\dist_T(y,v)
 \leq\alpha+1+\dist_T(y,v)=\dist_T(a,v).
\]
The pre-existing domination on $A$ is unchanged, so $a$ works.
Otherwise $\beta>\alpha+1$, which implies $\alpha\leq\beta+1$;
the same argument with the two sets exchanged shows that $b$ works.
\end{proof}

\begin{proposition}[Height-dominating root]\label{prop:upper-consolidation}
There exists a vertex $r\in V(T)$ such that
\begin{equation}\label{eq:upper-height-domination}
 h(v)\leq\dist_T(r,v)\qquad(v\in V(T)).
\end{equation}
\end{proposition}

\begin{proof}
Begin with the connected source partition of
Lemma~\ref{lem:upper-source-partition}; each part is height-dominated by
its own source.  If there is more than one part, connectedness of $T$
provides an edge joining two parts: for example, follow a path from a
vertex in one part to a vertex outside it and take its first exit edge.
Lemma~\ref{lem:upper-merge}, with $g=h$, merges those two parts and
supplies a root that height-dominates their union.  The resulting sets are
still a partition into connected, height-dominated parts.  Each merge
reduces the number of parts by one, so finitely many merges leave the
whole vertex set and one dominating root.
\end{proof}

\begin{theorem}[Arbitrary-defect upper bound]\label{thm:mass-bound}
For every configuration $C$ on a finite tree $T$ with at least two vertices,
\[
 \bigl(S_v(C)\leq0\text{ for all }v\bigr)
 \quad\Longrightarrow\quad
 \mass{C}\leq\estim(T)-1.
\]
\end{theorem}

\begin{proof}
For $C=0$, the bound follows from the positive estimator formula.  Otherwise
the preceding construction applies.  Choose the root $r$ from
Proposition~\ref{prop:upper-consolidation}.  Since $2^x$ increases with
$x$, Lemma~\ref{lem:upper-leaf-slack} gives
\begin{align*}
 \mass{C}
 &\leq L+\sum_{\deg_T(v)>1}\deg_T(v)2^{h(v)}\\
 &\leq L+\sum_{\deg_T(v)>1}\deg_T(v)2^{\dist_T(r,v)}\\
 &\leq L+\max_{s\in V(T)}
              \sum_{\deg_T(v)>1}\deg_T(v)2^{\dist_T(s,v)}
   =\estim(T)-1.
\end{align*}
The final equality is the estimator identity~\eqref{eq:compact-estimator}.
Every quantity in this calculation uses the ambient tree.
\end{proof}

Together with the exact score criterion (Theorem~\ref{thm:root-score}), Theorem~\ref{thm:mass-bound}
implies that every non-stackable configuration has mass at most
$\estim(T)-1$, and hence every configuration of mass $\estim(T)$ is
stackable.
